%==========================================================================
\section{Grammar extensions}
\begin{center}\small
\begin{tabularx}{\linewidth}{@{}l >{\raggedright}p{7.6cm} X@{}}
\toprule
\textbf{Id} & \textbf{Production} & \textbf{Comments} \\
\midrule
Added P54 & \nt{Expr} $\to$ \code{reflect} ( \nt{Expr} ) &
  Returns a ReflectInfo for the object's runtime class. \\
Added P55 & \nt{Expr} $\to$ \code{getfield} ( \nt{Type} , \nt{Expr} , \nt{Expr} ) &
  Reads the field whose name is a String value. \\
Added P56 & \nt{Expr} $\to$ \code{invoke} ( \nt{Type} , \nt{Expr} , \nt{Expr} (, \nt{Actuals})? ) &
  Calls a \textbf{non-void} method by name. Parallels P23. \\
Added P57 & \nt{Expr} $\to$ \code{new\_instance} ( \nt{Type} , \nt{Expr} (, \nt{Actuals})? ) &
  Creates an object of the class named by a String value; the class must be T
  or a subclass of T. \\
Added P58 & \nt{Stmt} $\to$ \code{setfield} ( \nt{Type} , \nt{Expr} , \nt{Expr} , \nt{Expr} ) ; &
  Writes the field whose name is a String value. \\
Added P59 & \nt{Stmt} $\to$ \code{invoke} ( \nt{Type} , \nt{Expr} , \nt{Expr} (, \nt{Actuals})? ) ; &
  Calls a \textbf{void} method by name. Parallels P19. \\
Added P60 & \nt{Expr} $\to$ \code{allocate} ( \nt{Type} , \nt{Expr} ) &
  Creates an object of the class named by a String value without running a
  constructor; the class must be T or a subclass of T. \\
\bottomrule
\end{tabularx}
\end{center}

No productions are modified or deleted. \code{reflect}, \code{getfield},
\code{setfield}, \code{invoke}, \code{new\_instance} and \code{allocate} become reserved
keywords; using one as an identifier is
\code{E\_RESERVED\_KEYWORD\_AS\_IDENTIFIER}. A new production begins with its
own keyword, and no LO-4 alternative of \nt{Expr} or \nt{Stmt} can begin with
that keyword, so the parser picks each production with one token of
lookahead. The grammar remains LL(2). As with P23 and P19, the parser doesn't
distinguish void from non-void invoke; the type checker does.

